\subsection{实现细节}
\label{sec:implementation}

本文基于 PyTorch 实现所提出的方法，并采用 CLIP ViT-B/32 作为图像和文本编码器。输入图像被调整为 $224\times224$，文本最大长度设置为 77 个 token。除非特别说明，检索分支采用双向 InfoNCE 损失，初始温度为 0.02。模型使用 Adam 优化器训练，学习率为 $1\times10^{-6}$，权重衰减为 $4\times10^{-5}$，动量参数为 $(0.9,0.999)$，batch size 为 64，并使用余弦学习率衰减和自动混合精度训练。

在噪声图文关联实验中，我们遵循常用的合成噪声协议，随机替换一定比例的文本描述以构造错配图文对。样本 clean/noisy 状态由两成分 GMM 根据逐样本对比损失估计。初始全数据 GMM 阈值设为 0.95，后续 loss-bank GMM 阈值设为 0.7。前 5 个 epoch 作为 warm-up 阶段，在该阶段中疑似噪声样本被抑制，以降低早期确认偏差。

clean memory 默认大小为 4096。只有置信度高于当前 clean 样本平均置信度并满足额外裕量的样本会被写入 memory。启用层级校准前，clean memory 至少需要包含 1024 个样本。层级结构默认由 clean memory 中的可信图文特征构建，融合系数设为 $\alpha=0.5$。每隔 100 个 training iteration 使用最多 4096 个 memory 样本刷新一次层级结构。

层级发现采用质量感知凝聚式聚合。首先基于融合特征、图像特征和文本特征共同计算跨模态亲和度，并为每个样本保留 $K_g=20$ 个近邻以构造稀疏亲和图。跨模态亲和度中的融合权重 $\beta$ 默认设为 0.5。初始微簇最大规模设为 16 或 32。凝聚式合并时同时考虑跨节点亲和度、节点一致性、跨模态可靠性和子节点分离度，并将得到的节点质量 $Q_k$ 存入树结构。低质量节点仍可作为层级路径的一部分，但不会触发强校准。

对于疑似噪声样本，图像特征和文本特征分别查询同一棵质量感知原型树，得到视觉叶节点和文本叶节点。每个叶节点到根节点的路径在建树后预先缓存，最深公共祖先可以通过比较两条路径的最长公共前缀得到。最终校准目标为该公共祖先节点的融合原型。校准权重同时由样本先验、祖先节点质量、查询置信度和祖先深度决定。低查询置信度样本、公共祖先过浅样本或祖先质量较低样本不会被强校准，而是使用较小权重进行弱校准。

双曲表示作为层级几何增强模块使用。默认双曲投影维度为 256，曲率为 1.0。层级正则包含父子紧致、深度顺序和兄弟分离三部分，并使用节点质量作为几何正则权重。为减少原始检索目标和祖先校准目标之间的梯度冲突，启用 HHNC 时疑似噪声样本在检索分支中的权重乘以 $\lambda_{ret}=0.5$，祖先校准损失权重设为 0.1。低置信弱校准权重设为 0.1。

主要命令行开关如下：
\begin{verbatim}
--enable_hhnc
--hhnc_hierarchy_source fused
--hhnc_fusion_alpha 0.5
--hhnc_loss_weight 0.1
--hhnc_low_conf_weight 0.1
--hhnc_noisy_ret_weight 0.5
--hhnc_hyper_dim 256
--hhnc_curvature 1.0
\end{verbatim}

为验证各个设计的贡献，建议进行如下消融实验：文本树与融合树对比，递归二均值树与质量感知凝聚树对比，HDBSCAN 初始化与微簇初始化对比，欧式原型树与双曲原型树对比，扁平原型校准与层级祖先校准对比，单路径校准与双路径 LCA 校准对比，固定树与在线刷新树对比，强制校准与质量感知弱校准对比，以及 retrieval-only correction 与 HHNC-enhanced correction 对比。
